\documentclass[12pt]{article}
\usepackage[T1]{fontenc}
\usepackage[utf8]{inputenc}
\usepackage{lmodern}
\usepackage{textcomp}
\usepackage{enumitem}
\usepackage{geometry}
\geometry{margin=1in}
\begin{document}
\begin{center}
Answer Key - Chemistry - Graduate Level Assessment 14 \
\small (Answers are ordered exactly as in the question paper.)
\end{center}

\section*{Multiple Choice}
\begin{enumerate}[label=\arabic*.]
\item \textbf{Answer:} A
\\ \textit{Options:} A) 5.4 moles per mile; B) 9.0 moles per mile; C) 2.5 moles per mile; D) 7.3 moles per mile; E) 3.8 moles per mile
\\ \textit{Note:} The correct answer is 5.4 moles per mile because the mass of 0.33 lb of CO can be converted to moles using the molar mass of CO (28.01 g/mol), resulting in approximately 5.4 moles of CO produced per mile.
\item \textbf{Answer:} A
\\ \textit{Options:} A) 350 lbs. sulfur; B) 320 lbs. sulfur; C) 425 lbs. sulfur; D) 400 lbs. sulfur; E) 300 lbs. sulfur
\\ \textit{Note:} The correct answer is 350 lbs. of sulfur because the stoichiometric ratio of aluminum to sulfur in aluminum sulfide (Al$_{2}$S$_{3}$) requires three moles of sulfur (96.06 g/mol) for every two moles of aluminum (26.98 g/mol), which translates to a mass ratio where 600 lbs of Al$_{2}$S$_{3}$ necessitates 350 lbs of sulfur to achieve the desired product.
\item \textbf{Answer:} D
\\ \textit{Options:} A) 8500 $\mathrm{~K}$; B) 15000 $\mathrm{~K}$; C) 9500 $\mathrm{~K}$; D) 11000 $\mathrm{~K}$
\\ \textit{Note:} The correct answer is 11000 K because using Wien's displacement law, which states that the peak wavelength of a blackbody spectrum is inversely proportional to its temperature, we can calculate the temperature of Sirius by rearranging the formula \( T = \frac{b}{\lambda_{\max}} \) where \( b \) is Wien's displacement constant and \( \lambda_{\max} \) is given as 260 nm.
\item \textbf{Answer:} A
\\ \textit{Options:} A) 43 minutes; B) 103 minutes; C) 93 minutes; D) 113 minutes; E) 23 minutes
\\ \textit{Note:} The correct answer is 43 minutes because the reaction follows first-order kinetics, where the time required for a substance's concentration to decrease to a certain percentage can be calculated using the half-life, leading to the conclusion that it takes approximately 43 minutes to reduce the concentration of SO\_2 Cl\_2 to 1.0\% of its original amount.
\item \textbf{Answer:} A
\\ \textit{Options:} A) C\_3 H\_6; B) C\_5 H\_10; C) C\_4 H\_8; D) C\_4 H\_9; E) C\_3 H\_8
\\ \textit{Note:} The molecular formula C$_{3}$H$_{6}$ corresponds to a compound with a molecular weight of 58 g/mol, which matches the given percentage composition of 82.8\% carbon and 17.2\% hydrogen when calculated based on the molar mass of each element.
\end{enumerate}

\section*{True / False}
\begin{enumerate}[label=\arabic*.]
\setcounter{enumi}{5}
\item \textbf{Answer:} True
\\ \textit{Options:} A) True; B) False
\\ \textit{Note:} The final concentration of the mixed solution is 0.35 M H₂SO₄ because the total moles of H₂SO₄ from both solutions combined (25 mmol from the first solution and 18.75 mmol from the second solution) divided by the total volume (125 ml) results in a concentration of 0.35 M.
\item \textbf{Answer:} False
\\ \textit{Options:} A) True; B) False
\\ \textit{Note:} The statement is false because the number of configuration state functions generated in a full configuration interaction calculation with a 6-31 G** basis set for CH 2 SiHF is significantly lower than 4.67×$10^28$ due to the limitations in the number of active electrons and orbitals considered in the calculation.
\item \textbf{Answer:} False
\\ \textit{Options:} A) True; B) False
\\ \textit{Note:} The statement is false because the stoichiometric ratio between NH 3 and O$_{2}$ in the balanced chemical equation does not require 72 g of NH 3 to fully react with 80 g of O$_{2}$, indicating that the quantities of reactants must be calculated based on their molar ratios.
\item \textbf{Answer:} True
\\ \textit{Options:} A) True; B) False
\\ \textit{Note:} The statement is true because, according to the Heisenberg Uncertainty Principle, the uncertainty in momentum is inversely related to the uncertainty in position, and a position uncertainty of 10 pm leads to a calculated momentum uncertainty of approximately 4.2 x 10^-23 kg·m·s^-1 for an electron.
\item \textbf{Answer:} True
\\ \textit{Options:} A) True; B) False
\\ \textit{Note:} During boiling, the thermal energy supplied to liquid NH 3 is sufficient to overcome the hydrogen bonds between NH 3 molecules, causing them to break and allowing the molecules to enter the gas phase.
\end{enumerate}

\section*{Long Form Response}
\begin{enumerate}[label=\arabic*.]
\setcounter{enumi}{10}
\item \textbf{Answer:} The relationship between kinetic energy and wavelength is articulated through the de Broglie wavelength formula, which describes the wave-particle duality of matter. For a photon with 1 eV of kinetic energy, its wavelength can be calculated using the equation λ = hc/E, resulting in a wavelength of approximately 1300 nm. In contrast, for an electron possessing the same kinetic energy, the de Broglie wavelength is given by λ = h/p, where p is momentum; this yields a wavelength of approximately 1.300 nm. The significant difference in wavelengths-1300 nm for the photon versus 1.300 nm for the electron- highlights the contrasting wave behaviors of massless and massive particles, indicating that photons exhibit longer wavelengths due to their lack of mass, while electrons, possessing mass, exhibit much shorter wavelengths, emphasizing the implications of their respective wave-particle dualities in quantum mechanics.
\\ \textit{Note:} The answer correctly identifies the application of the de Broglie wavelength formula to relate kinetic energy to wavelength for both a photon and an electron, demonstrating that while the photon has a significantly longer wavelength due to its massless nature, the electron exhibits a much shorter wavelength reflecting its particle characteristics and mass.
\item \textbf{Answer:} At standard temperature and pressure (STP), which is defined as 0 degrees Celsius and 1 atmosphere of pressure, the density of air is approximately 1.29 grams per liter. To find the weight of 1,000 cubic feet of air, we first convert cubic feet to liters, knowing that 1 cubic foot equals approximately 28.3168 liters. Thus, 1,000 cubic feet is about 28,316.8 liters. Multiplying the volume in liters by the density gives us a mass of approximately 36,519 grams. Converting grams to ounces (1 ounce = 28.3495 grams), we find the weight to be approximately 1.10 \ensuremath{\times} $10^3$ ounces. This calculation illustrates how gas density and volume interplay to determine the mass of air at STP conditions.
\\ \textit{Note:} correct because it accurately applies the standard density of air at STP, converts the volume from cubic feet to liters, and uses appropriate unit conversions to calculate the weight of the air.
\end{enumerate}

\vfill
\noindent\textit{End of Answer Key}
\end{document}